% Context from chapters/convex-analysis.tex; for mathematical reference, not compilation.
\section{Basics of Convex Analysis}

We consider minimization problems of the form
\eql{\label{eq-convex-function}
	\umin{x \in \Hh} f(x)
}
over the finite-dimensional Hilbert space $\Hh \eqdef \RR^N$, with the canonical inner product $\dotp{\cdot}{\cdot}$.
 
Many of these results extend to infinite-dimensional Hilbert spaces, with additional hypotheses where compactness is used.

The objective $f : \Hh \longrightarrow \bar \RR \eqdef \RR \cup \{+\infty\}$ is convex. Allowing $f$ to take the value $+\infty$ encodes constraints directly in the objective; its effective domain is
\eq{
	\dom(f) \eqdef \enscond{x}{f(x)<+\infty}.
}

For a set $\Cc \subset \Hh$, define the indicator function
\eq{
	\iota_{\Cc}(x) \eqdef \choice{
		0 \qifq x \in \Cc, \\
		+\infty \quad\text{otherwise.}
	}
}


                                                                                  
\subsection{Convex Sets and Functions}

A set $\Om \subset \Hh$ is convex if
\eq{
	\foralls (x,y,t) \in \Om^2 \times [0,1], \quad (1-t)x+ty \in \Om.
}
A function is convex if
\eql{\label{eq-convex-defn}
	\foralls (x,y,t) \in \Hh^2 \times [0,1], \quad f( (1-t)x+ty ) \leq (1-t)f(x) + t f(y)
}
Equivalently, its epigraph $\enscond{(x,r) \in \Hh \times \RR}{r \geq f(x)}$ is convex. The inequality is interpreted in the extended real line $\bar\RR$, with the convention $0\cdot(+\infty)=0$ at the endpoints $t=0,1$.
 
The function $f$ is strictly convex if the inequality in~\eqref{eq-convex-defn} is strict whenever $x,y\in\dom(f)$ are distinct and $0<t<1$.
 
A set $\Om$ is convex if and only if $\iota_{\Om}$ is a convex function.



\begin{figure}[tbp]
\centering
\BookDrawing{tikz/convex-analysis-convex-examples}
\caption{Convexity and strict convexity for functions and sets.}
\label{fig-convex-examples}
\end{figure}




Throughout the chapter, convex functions $f$ are assumed proper, meaning $\dom(f) \neq \emptyset$, and lower semicontinuous (lsc), meaning that for every $x \in \Hh$,
\eq{
	\liminf_{y \rightarrow x} f(y)  \geq f(x).
}
For a convex function, lower semicontinuity is equivalent to the epigraph $\epi(f)$ being closed.
 
We denote by $\Ga_0(\Hh)$ the set of proper convex lsc functions.

                                                                                  
\subsection{First-Order Conditions}

   
\paragraph{Existence of minimizers.}

We first address existence. Lower semicontinuity and coercivity provide a useful sufficient condition.

\begin{prop}
	If $f$ is proper, lower semicontinuous, and coercive (i.e. $f(x)\rightarrow +\infty$ as $\norm{x}\rightarrow +\infty$), then there exists a minimizer $x^\star$ of $f$.
\end{prop}

\begin{proof}
Choose $x_0\in\dom(f)$. The sublevel set $K=\{x:f(x)\leq f(x_0)\}$ is nonempty, closed by lower semicontinuity, and bounded by coercivity. It is therefore compact in finite dimensions. A lower semicontinuous function attains its infimum on $K$, at some $x^\star$. Since $f(x)>f(x_0)\geq f(x^\star)$ outside $K$, this point is a global minimizer.
\end{proof}

This compactness argument is known as the direct method in the calculus of variations.
 
For $f$ in $\Ga_0(\Hh)$, the minimizer set $\argmin f$ is closed and convex, and every local minimizer is global. Strict convexity of $f$ implies that there is at most one minimizer.

   
\paragraph{Subdifferential.}

The subdifferential at $x$ of $f$ is
\eq{
	\partial f(x) \eqdef \enscond{u \in \Hh^*}{ \forall y\in\Hh, f(y) \geq f(x) + \dotp{u}{y-x} },\qquad x\in\dom(f).
}
We set $\partial f(x)=\emptyset$ outside $\dom(f)$. Here $\Hh^*=\RR^N$ denotes the space of dual vectors. The inner product identifies the dual space with $\Hh$, but distinguishing primal and dual variables remains useful. The duality pairing itself is canonical; the identification of a linear functional with 